\unchapter{摘\quad 要}

无人机网络可快速部署于灾后应急观测、低空巡检和临时通信覆盖场景，完成信息采集与协同判断。节点移动、链路波动和有限能量可能造成中继负载集中与回传中断；不同视角的语义证据价值不一，无差别交换会增加开销；数据丢失和确认不完整又使已选证据难以及时交付。本文针对持续传输、内容选择和可靠交付中的这些问题，研究资源受限无人机网络的通信优化方法。

（1）针对高动态网络中的中继能量热点，提出基于李雅普诺夫优化的能量感知分布式路由方法 LEADR。建立节点移动、无线链路、数据队列和能量变化模型，以虚拟能量队列记录长期支出压力，将跨时隙问题转化为局部在线决策。节点按链路条件求解功率，结合邻居剩余能量和地理进度选择下一跳，并依据相对运动预测调整邻居更新频率。节点数量、通信范围和飞行速度实验表明，该方法在维持任务服务周期的同时保持了较高转发吞吐与较低丢包率。

（2）针对证据价值、消息粒度与链路成本的耦合，提出链路感知选择性语义通信方法 EviComRL。冻结视觉语言模型生成五级证据可见性轨迹，语义提取与策略学习分别进行。策略联合编码发送方、中心需求、粒度、链路、可靠性及已选历史，通过带动作掩码的图结构网络和近端策略优化选择消息与停止时机。主设置下准确率为58.17\%，宏平均 F1 为50.58\%；相较全通信与上下文老虎机，准确率提高11.00和7.50个百分点。与多层感知机策略相比，分类质量提高且每任务平均通信量减少。多条件实验分析了覆盖、预算、消息表达和可靠性特征的影响。

（3）针对未确认发送造成的接收不确定性与重复通信，提出一致性感知的信念反馈控制方法 C-BFC。轻量网络估计接收集合效用，同类别一致性投影修正同任务、同类别非空集合之间的虚假增量。控制器依据累计确认和发送历史递推接收分布，以两步发送前瞻和反馈等待价值比较新发、重传、等待与停止，并通过跨任务虚拟能量压力调节发送代价。在17个主比较条件下，平均按时正确完成率为49.84\%，单位正确任务通信能耗为3.615 mJ。相较同环境重新训练的 SD-SAC 和 TP-DQN，正确完成率提高0.87和0.66个百分点，单位正确任务能耗降低16.59\%和14.32\%。

\vspace{1em}
\noindent\textbf{关键词：}无人机网络；能量感知路由；选择性语义通信；强化学习；不完备反馈

\unchapter{Abstract}

UAV networks can be deployed rapidly for post-disaster observation, low-altitude inspection, and temporary communication coverage. Mobility, fluctuating links, and limited energy can concentrate relay loads and interrupt data delivery. Local observations provide evidence of unequal value, and indiscriminate exchange increases communication costs. Data loss and incomplete acknowledgments can also prevent timely delivery of selected evidence. This thesis studies these problems at three stages of communication in resource-constrained UAV networks: sustained transmission, semantic selection, and reliable delivery.

First, LEADR is proposed as a Lyapunov-based energy-aware decentralized routing method for dynamic UAV networks. Its model includes mobility, wireless links, data queues, and energy evolution. Virtual energy queues record accumulated expenditure pressure and convert allocation across time slots into local online decisions. Nodes determine transmission power for candidate links and select the next hop using neighbor residual energy and geographic progress. Relative-motion prediction sets the frequency of neighbor-state updates. Experiments varying UAV count, communication range, and flight speed show that LEADR sustains service over the task period while maintaining high forwarding throughput and low packet loss.

Second, EviComRL selects semantic messages according to link conditions and task needs. Frozen vision-language models generate structured traces for a five-level evidence-visibility task, so semantic extraction and policy learning can be performed separately. Candidate representations encode sender information, receiver demand, message granularity, link conditions, reliability, and selected-message history. A graph-structured policy with action masking is trained through proximal policy optimization to select messages and determine when to stop. In the main setting, EviComRL achieves 58.17\% accuracy and 50.58\% macro F1. Accuracy exceeds full communication and a contextual-bandit controller by 11.00 and 7.50 percentage points, respectively. Compared with an MLP policy, it improves classification quality while reducing average token expenditure per task. Additional experiments examine coverage, budgets, message modes, and reliability features.

Third, consistency-aware belief feedback control, C-BFC, is developed for uncertain reception and repeated transmissions caused by missing acknowledgments. A lightweight network predicts received-set utility. A class-consistency projection removes spurious estimated gains between nonempty sets with the same predicted class for the same task. Cumulative acknowledgments and transmission history update the reception-state distribution. Two-step transmission lookahead and the value of waiting for feedback guide new transmissions, retransmissions, waiting, and stopping. Virtual energy pressure adjusts transmission costs across tasks. Across seventeen main comparison conditions, C-BFC achieves a mean on-time correct completion rate of 49.84\% and uses 3.615 mJ per correctly completed task. Relative to SD-SAC and TP-DQN retrained in the same environment, completion rates increase by 0.87 and 0.66 percentage points, while energy per correct task decreases by 16.59\% and 14.32\%, respectively.

\vspace{1em}
\noindent\textbf{Keywords:} UAV networks; energy-aware routing; selective semantic communication; reinforcement learning; incomplete feedback
